\subsection{Conjuntos}

\begin{itemize}
    \item $C$: Conjunto de locais candidatos para instalação de eletropostos
    \item $P$: Conjunto de Pontos de Interesse (POIs) -- hospitais, educação, transporte, comércio, recreação
    \item $V(c)$: Conjunto de vizinhos espaciais do candidato $c$ (dentro de raio crítico de influência)
    \item $\mathcal{S}_c$: Conjunto de todas as combinações possíveis de vizinhos de $c$ que podem ser designados simultaneamente ($\mathcal{S}_c \subseteq 2^{V(c)}$)
\end{itemize}

\subsection{Parâmetros}

\subsubsection{Lucro (pré-computado via simulador)}

\begin{itemize}
    \item $R_c^{\text{próprio}}$: Lucro líquido do eletroposto $c$ ao longo do horizonte de 10 anos, se operasse isoladamente (R\$). Computado uma única vez pelo simulador em etapa de pré-processamento, considerando custos de instalação, operação e receita de atendimento de demanda.
    
    \item $\Delta R_{c,c'}$: Lucro desviado do candidato $c$ para seu vizinho $c' \in V(c)$ quando $c'$ é designado. Representa a parcela de demanda que $c$ perde quando $c'$ entra em operação (R\$). Também pré-computado, contemplando:
    \begin{itemize}
        \item Distância real pela malha viária (não euclidiana)
        \item Preferência de usuários (modelo de escolha discreta)
        \item Viabilidade de rede elétrica (se ambos $c$ e $c'$ podem operar simultaneamente em hora pico)
    \end{itemize}
    
    \item $R_c(S_c)$: Lucro esperado do candidato $c$ quando exatamente o conjunto de vizinhos $S_c \subseteq V(c)$ está designado. Definido como:
    \begin{equation}
    R_c(S_c) = R_c^{\text{próprio}} - \sum_{c' \in S_c} \Delta R_{c,c'}
    \end{equation}
    Todos estes valores são pré-computados na simulação para todas as $2^{|V(c)|}$ combinações possíveis.
\end{itemize}

\subsubsection{Qualidade de rede}

\begin{itemize}
    \item $Q_{\min}^{\text{cob}}$: Cobertura mínima de POIs requerida (percentual, ex: 85\%)
    \item $r_{\text{cob}}$: Raio de cobertura máximo aceitável para um POI ser considerado atendido (metros)
\end{itemize}

\subsection{Variáveis de decisão}

\subsubsection{Variáveis principais}

\begin{itemize}
    \item $x_c \in \{0,1\}$: Indicadora -- designar eletroposto no candidato $c$
    \item $y_p \in \{0,1\}$: Indicadora -- POI $p$ está coberto por pelo menos um eletroposto designado
\end{itemize}

\subsubsection{Variáveis de configuração (para linearização)}

\begin{itemize}
    \item $z_{c,S_c} \in \{0,1\}$: Indicadora -- a configuração exata de vizinhos designados para $c$ é o conjunto $S_c$. \textbf{Nota:} Esta variável é opcional em implementação -- pode ser eliminada se o solver suportar formulação direta com tabela de retornos pré-computados (recomendado para eficiência computacional).
\end{itemize}

\subsection{Função objetivo}

\subsubsection{Versão Bilineal (Conceitual -- Não Recomendada para Resolver)}

\begin{equation}
\label{eq:objetivo_bilineal}
\max \quad Z = \sum_{c \in C} x_c \cdot \left( R_c^{\text{próprio}} - \sum_{c' \in V(c)} x_{c'} \cdot \Delta R_{c,c'} \right)
\end{equation}

\textit{Problema:} Esta formulação é \textbf{bilineal} -- os termos $x_c \cdot x_{c'}$ a tornam não-linear e intratável para solvers MILP estândar.

\subsubsection{Versão Linear (Recomendada -- Usando Pré-cómputo)}

\begin{equation}
\label{eq:objetivo_linear}
\max \quad Z = \sum_{c \in C} \sum_{S_c \in \mathcal{S}_c} R_c(S_c) \cdot z_{c,S_c}
\end{equation}

\textit{onde} $R_c(S_c)$ \textit{é o lucro pré-computado do candidato } $c$ \textit{ quando exatamente o conjunto de vizinhos } $S_c$ \textit{ está designado.}

\vspace{0.5cm}

\textbf{Vantagem:} O modelo é agora \textbf{linear puro} -- altamente eficiente para solvers MILP.

\subsection{Restrições}

\subsubsection{Grupo 1: Configuração de Vizinhos}

\hspace{0.5cm} \textbf{R1: Atribuição de configuração}

\begin{align}
\sum_{S_c \in \mathcal{S}_c} z_{c,S_c} &= x_c \quad \forall c \in C \label{eq:restricao_config}
\end{align}

Um candidato $c$ designado deve ter exatamente uma configuração de vizinhos ativa.

\hspace{0.5cm} \textbf{R2: Consistência de designações com configurações}

Para cada candidato $c$ e cada configuração $S_c$, se $z_{c,S_c} = 1$, então todos os vizinhos em $S_c$ devem estar designados:

\begin{align}
x_{c'} &\geq z_{c,S_c} \quad \forall c \in C, c' \in S_c, S_c \in \mathcal{S}_c \label{eq:restricao_consistencia}
\end{align}

Equivalentemente, para vizinhos \textit{não} em $S_c$:

\begin{align}
x_{c'} &\leq 1 - z_{c,S_c} \quad \forall c \in C, c' \in V(c) \setminus S_c, S_c \in \mathcal{S}_c \label{eq:restricao_exclusao}
\end{align}

\subsubsection{Grupo 2: Qualidade de Rede}

\hspace{0.5cm} \textbf{R3: Indicador de cobertura de POI}

\begin{align}
y_p &\leq \sum_{c: \text{dist}(c,p) \leq r_{\text{cob}}} x_c \quad \forall p \in P \label{eq:cobertura_poi}
\end{align}

Um POI $p$ é considerado coberto se existe pelo menos um eletroposto designado dentro do raio máximo aceitável.

\hspace{0.5cm} \textbf{R4: Meta de cobertura mínima}

\begin{align}
\sum_{p \in P} y_p &\geq Q_{\min}^{\text{cob}} \cdot |P| \label{eq:meta_cobertura}
\end{align}

A quantidade total de POIs cobertos deve atender ao mínimo especificado.

\subsection{Algoritmo de Solução}

O modelo é resolvido em uma única passada usando um solver MILP estândar (CPLEX, Gurobi):

\begin{enumerate}
    \item \textbf{Pré-cómputo}: O simulador calcula $R_c(S_c)$ para todos os candidatos e todas as $2^{|V(c)|}$ combinações de vizinhos
    \item \textbf{Resolução}: O solver MILP maximiza a função objetivo linear (Equação \ref{eq:objetivo_linear}), satisfazendo R1--R4
    \item \textbf{Pós-processamento}: Extrair a solução ótima $x_c^*$ e identificar qual configuração de vizinhos foi selecionada para cada $c$
\end{enumerate}

\textbf{Complexidade computacional:} Para $|C|$ candidatos e promédio de $\bar{d}$ vizinhos por candidato:
\begin{itemize}
    \item \textbf{Variáveis binárias}: $O(|C| \cdot 2^{\bar{d}})$ (controlável se $\bar{d} \leq 10$)
    \item \textbf{Restrições}: $O(|C| \cdot 2^{\bar{d}} + |P|)$
    \item \textbf{Tempo de resolução esperado}: Segundos a minutos para instâncias moderadas
\end{itemize}

\subsection{Justificativa de Simplicidade e Eficiência}

Esta abordagem é superior porque:

\begin{itemize}
    \item \textbf{Totalmente linear}: Sem termos bilineais, integráveis em solvers MILP estândar
    \item \textbf{Sem ajustes pós-hoc}: Toda factibilidade (rede elétrica, demanda) está no pré-cómputo
    \item \textbf{Transparência de lucro}: Objetivo expressa claramente qual lucro se obtém por cada combinação de designações
    \item \textbf{Escalabilidade}: Se o número de vizinhos por candidato é pequeno ($\leq 10$), o modelo é tratável
    \item \textbf{Separação de responsabilidades}: Simulador (simulação física) vs. Modelo (otimização discreta)
\end{itemize}

\subsection{Extensões Futuras}

Se em fases posteriores for necessário:

\begin{itemize}
    \item \textbf{Granularidade temporal}: Discretizar o horizonte em períodos e ter $R_c(S_c, t)$ por período
    \item \textbf{Restrições adicionais}: Orçamento global $\sum_c \text{CAPEX}_c \cdot x_c \leq B$
    \item \textbf{Múltiplos critérios}: Incorporar critérios ambientais/sociais via ponderação no objetivo
    \item \textbf{Meta-heurísticas}: Implementar algoritmos genéticos ou busca tabu para comparação
    \item \textbf{Robustez}: Otimização robusta contra incerteza em parâmetros de lucro
\end{itemize}

O modelo base permanece linear e manejável, permitindo evolução incremental.

\subsection{Bounds Teóricos}

O lucro máximo teórico (sem restrições de cobertura):

\begin{equation}
Z_{\max} = \sum_{c \in C} R_c(\emptyset)
\end{equation}

onde $R_c(\emptyset)$ é o lucro do candidato $c$ se nenhum vizinho fosse designado.

O mínimo (nenhum eletroposto designado):

\begin{equation}
Z_{\min} = 0
\end{equation}

Com restrições de qualidade de rede, o lucro será reduzido conforme a necessidade de garantir cobertura mínima de POIs force a designação de candidatos com menor lucro marginal.